\documentclass[11pt]{article}
\usepackage{amsmath, amssymb, geometry, graphicx, enumerate}
\geometry{a4paper, margin=1in}

\title{\textbf{Homework 1.3: Vector Equations}}
\author{Tashfeen Omran}
\date{September 2026}

\begin{document}

\maketitle
\tableofcontents
\vspace{1cm}
\hrule
\vspace{1cm}

\section{Part 1: Comprehensive Introduction, Context, and Prerequisites}

\subsection{Core Concepts}
In the previous sections, we treated systems of linear equations simply as arrays of numbers. In Section 1.3, we introduce a powerful new geometric and algebraic perspective by organizing these numbers into \textbf{vectors}. 

A \textbf{column vector} (or simply a vector) is a matrix with only one column. When a vector has $n$ entries, we say it belongs to $\mathbb{R}^n$ (read "R-n"), which represents the set of all $n$-tuples of real numbers.
For example, $\mathbf{u} = \begin{bmatrix} 3 \\ -1 \end{bmatrix}$ is a vector in $\mathbb{R}^2$, and it geometrically represents a directed line segment (an arrow) from the origin $(0,0)$ to the point $(3, -1)$ in a 2D coordinate plane.

We define two fundamental operations on vectors:
\begin{enumerate}
    \item \textbf{Vector Addition:} Two vectors of the same size are added entry-by-entry. Geometrically, this follows the "parallelogram rule" or "tip-to-tail" method.
    \item \textbf{Scalar Multiplication:} A vector can be multiplied by a real number (a scalar), which scales every entry in the vector by that number. Geometrically, this stretches, shrinks, or reverses the direction of the vector.
\end{enumerate}

\subsection{Linear Combinations and Span}
The true power of this section comes from combining these two operations. Given a set of vectors $\mathbf{v}_1, \mathbf{v}_2, \dots, \mathbf{v}_p$ and scalars $c_1, c_2, \dots, c_p$, the expression:
\[ \mathbf{y} = c_1\mathbf{v}_1 + c_2\mathbf{v}_2 + \dots + c_p\mathbf{v}_p \]
is called a \textbf{linear combination} of the vectors with weights $c_1, \dots, c_p$.

The set of \textit{all possible} linear combinations of a set of vectors is called the \textbf{Span} of those vectors, denoted as $\text{Span}\{\mathbf{v}_1, \dots, \mathbf{v}_p\}$. 
\begin{itemize}
    \item The Span of a single non-zero vector $\text{Span}\{\mathbf{v}\}$ is geometrically a \textbf{line} through the origin.
    \item The Span of two non-collinear vectors $\text{Span}\{\mathbf{u}, \mathbf{v}\}$ is geometrically a \textbf{plane} through the origin.
\end{itemize}

\subsection{Intuition and Derivation: The Fundamental Equivalence}
Why do we care about linear combinations? Because asking "Can vector $\mathbf{b}$ be generated by a linear combination of the columns of matrix $A$?" is mathematically identical to asking "Does the linear system corresponding to the augmented matrix $[A \ \mathbf{b}]$ have a solution?" 
The vector equation $x_1\mathbf{a}_1 + x_2\mathbf{a}_2 + \dots + x_n\mathbf{a}_n = \mathbf{b}$ has the exact same solution set as the linear system whose augmented matrix is $[\mathbf{a}_1 \ \mathbf{a}_2 \ \dots \ \mathbf{a}_n \ \mathbf{b}]$. This bridges the gap between abstract geometry and concrete algebra.

\subsection{Historical Context and Motivation}
The concept of vectors emerged in the late 19th century through the work of physicists and mathematicians like Josiah Willard Gibbs and Oliver Heaviside. They needed a mathematical framework to model physical quantities that had both magnitude and direction (like velocity, force, and magnetic fields). By treating columns of a matrix as geometric objects, modern mathematicians realized they could visualize complex algebraic systems, birthing the field of linear algebra.

\subsection{Prerequisites}
You must be fluent in translating a system of equations into an augmented matrix and performing Gaussian elimination (row reduction) to find pivots, as mastered in Sections 1.1 and 1.2.

\newpage
\section{Part 2: Detailed Homework Solutions}

\subsection*{Problem 1.3.1}
\textbf{Question:} Compute $\mathbf{u} + \mathbf{v}$ and $\mathbf{u} - 2\mathbf{v}$.
\[ \mathbf{u} = \begin{bmatrix} -5 \\ 6 \end{bmatrix}, \ \mathbf{v} = \begin{bmatrix} -4 \\ 3 \end{bmatrix} \]

\textbf{Solution:}
To add vectors, add their corresponding entries:
\[ \mathbf{u} + \mathbf{v} = \begin{bmatrix} -5 \\ 6 \end{bmatrix} + \begin{bmatrix} -4 \\ 3 \end{bmatrix} = \begin{bmatrix} -5 + (-4) \\ 6 + 3 \end{bmatrix} = \begin{bmatrix} -9 \\ 9 \end{bmatrix} \]
To compute $\mathbf{u} - 2\mathbf{v}$, first perform the scalar multiplication on $\mathbf{v}$, then subtract:
\[ \mathbf{u} - 2\mathbf{v} = \begin{bmatrix} -5 \\ 6 \end{bmatrix} - 2\begin{bmatrix} -4 \\ 3 \end{bmatrix} = \begin{bmatrix} -5 \\ 6 \end{bmatrix} - \begin{bmatrix} -8 \\ 6 \end{bmatrix} = \begin{bmatrix} -5 - (-8) \\ 6 - 6 \end{bmatrix} = \begin{bmatrix} 3 \\ 0 \end{bmatrix} \]
\\
\textbf{Answer:} 
Part 1: $\begin{bmatrix} -9 \\ 9 \end{bmatrix}$ \quad Part 2: $\begin{bmatrix} 3 \\ 0 \end{bmatrix}$


\subsection*{Problem 1.3.5}
\textbf{Question:} Write a system of equations that is equivalent to the given vector equation.
\[ x_1 \begin{bmatrix} 5 \\ -4 \\ 7 \end{bmatrix} + x_2 \begin{bmatrix} 3 \\ 0 \\ -8 \end{bmatrix} = \begin{bmatrix} 4 \\ -2 \\ 9 \end{bmatrix} \]

\textbf{Solution:}
Multiply the scalars into the vectors:
\[ \begin{bmatrix} 5x_1 \\ -4x_1 \\ 7x_1 \end{bmatrix} + \begin{bmatrix} 3x_2 \\ 0x_2 \\ -8x_2 \end{bmatrix} = \begin{bmatrix} 4 \\ -2 \\ 9 \end{bmatrix} \]
Add the vectors on the left side:
\[ \begin{bmatrix} 5x_1 + 3x_2 \\ -4x_1 + 0 \\ 7x_1 - 8x_2 \end{bmatrix} = \begin{bmatrix} 4 \\ -2 \\ 9 \end{bmatrix} \]
Two vectors are equal if and only if their corresponding entries are equal. This gives us the system:
\begin{align*}
    5x_1 + 3x_2 &= 4 \\
    -4x_1 \phantom{+ 0x_2} &= -2 \\
    7x_1 - 8x_2 &= 9
\end{align*}
\\
\textbf{Answer:} D.


\subsection*{Problem 1.3.6}
\textbf{Question:} Write a system of equations that is equivalent to the given vector equation.
\[ x_1 \begin{bmatrix} 3 \\ -2 \end{bmatrix} + x_2 \begin{bmatrix} 9 \\ 2 \end{bmatrix} + x_3 \begin{bmatrix} -2 \\ 1 \end{bmatrix} = \begin{bmatrix} 0 \\ 0 \end{bmatrix} \]

\textbf{Solution:}
Applying the scalar multiplication and vector addition rules across the rows:
Row 1: $3x_1 + 9x_2 - 2x_3 = 0$
Row 2: $-2x_1 + 2x_2 + 1x_3 = 0$
\\
\textbf{Answer:} C.


\subsection*{Problem 1.3.9}
\textbf{Question:} Write a vector equation that is equivalent to the given system of equations.
\begin{align*}
    \phantom{3x_1 + } x_2 + 4x_3 &= 0 \\
    3x_1 + 6x_2 - x_3 &= 0 \\
    -x_1 + 3x_2 - 7x_3 &= 0
\end{align*}

\textbf{Solution:}
We isolate the variables $x_1, x_2,$ and $x_3$ as scalar weights multiplying column vectors formed by the coefficients in the system.
The $x_1$ coefficients are $0, 3, -1$.
The $x_2$ coefficients are $1, 6, 3$.
The $x_3$ coefficients are $4, -1, -7$.
The constants on the right side are $0, 0, 0$.
Forming the linear combination:
\[ x_1 \begin{bmatrix} 0 \\ 3 \\ -1 \end{bmatrix} + x_2 \begin{bmatrix} 1 \\ 6 \\ 3 \end{bmatrix} + x_3 \begin{bmatrix} 4 \\ -1 \\ -7 \end{bmatrix} = \begin{bmatrix} 0 \\ 0 \\ 0 \end{bmatrix} \]
\\
\textbf{Answer:} $x_1 \begin{bmatrix} 0 \\ 3 \\ -1 \end{bmatrix} + x_2 \begin{bmatrix} 1 \\ 6 \\ 3 \end{bmatrix} + x_3 \begin{bmatrix} 4 \\ -1 \\ -7 \end{bmatrix} = \begin{bmatrix} 0 \\ 0 \\ 0 \end{bmatrix}$


\subsection*{Problem 1.3.11}
\textbf{Question:} Determine if $\mathbf{b}$ is a linear combination of $\mathbf{a}_1, \mathbf{a}_2, \mathbf{a}_3$.
\[ \mathbf{a}_1 = \begin{bmatrix} 1 \\ -2 \\ 0 \end{bmatrix}, \ \mathbf{a}_2 = \begin{bmatrix} 0 \\ 1 \\ 4 \end{bmatrix}, \ \mathbf{a}_3 = \begin{bmatrix} 4 \\ -6 \\ 8 \end{bmatrix}, \ \mathbf{b} = \begin{bmatrix} 4 \\ -1 \\ 28 \end{bmatrix} \]

\textbf{Solution:}
This is equivalent to determining if the augmented matrix $[\mathbf{a}_1 \ \mathbf{a}_2 \ \mathbf{a}_3 \ \mathbf{b}]$ is consistent.
\[ \begin{bmatrix} 1 & 0 & 4 & 4 \\ -2 & 1 & -6 & -1 \\ 0 & 4 & 8 & 28 \end{bmatrix} \]
Row reduction steps:
\[ \underrightarrow{R_2 \leftarrow R_2 + 2R_1} \begin{bmatrix} 1 & 0 & 4 & 4 \\ 0 & 1 & 2 & 7 \\ 0 & 4 & 8 & 28 \end{bmatrix} \]
\[ \underrightarrow{R_3 \leftarrow R_3 - 4R_2} \begin{bmatrix} 1 & 0 & 4 & 4 \\ 0 & 1 & 2 & 7 \\ 0 & 0 & 0 & 0 \end{bmatrix} \]
The matrix is now in echelon form. The pivots are in the first column (entry $(1,1)$) and the second column (entry $(2,2)$). There is no pivot in the augmented rightmost column, meaning the system is consistent. Thus, $\mathbf{b}$ \textit{is} a linear combination of the vectors.
\\
\textbf{Answer:} A. Vector $\mathbf{b}$ is a linear combination of $\mathbf{a}_1, \mathbf{a}_2,$ and $\mathbf{a}_3$. The pivots in the corresponding echelon matrix are in the first entry in the first column and the second entry in the second column.


\subsection*{Problem 1.3.13}
\textbf{Question:} Determine if $\mathbf{b}$ is a linear combination of the vectors formed from the columns of the matrix $A$.
\[ A = \begin{bmatrix} 1 & -4 & 2 \\ 0 & 3 & 5 \\ -4 & 16 & -8 \end{bmatrix}, \ \mathbf{b} = \begin{bmatrix} 3 \\ -6 \\ -3 \end{bmatrix} \]

\textbf{Solution:}
Set up the augmented matrix:
\[ \begin{bmatrix} 1 & -4 & 2 & 3 \\ 0 & 3 & 5 & -6 \\ -4 & 16 & -8 & -3 \end{bmatrix} \]
Eliminate the $-4$ in the third row:
\[ \underrightarrow{R_3 \leftarrow R_3 + 4R_1} \begin{bmatrix} 1 & -4 & 2 & 3 \\ 0 & 3 & 5 & -6 \\ 0 & 0 & 0 & 9 \end{bmatrix} \]
The last row translates to the equation $0 = 9$, which is a contradiction. The system is inconsistent. Therefore, there are no weights $x_1, x_2, x_3$ that can generate $\mathbf{b}$ from the columns of $A$.
\\
\textbf{Answer:} B. Vector $\mathbf{b}$ is not a linear combination of the vectors formed from the columns of the matrix $A$.


\subsection*{Problem 1.3.14}
\textbf{Question:} Determine if $\mathbf{b}$ is a linear combination of the vectors formed from the columns of the matrix $A$.
\[ A = \begin{bmatrix} 1 & -6 & -5 \\ 0 & 3 & 2 \\ 2 & -12 & 9 \end{bmatrix}, \ \mathbf{b} = \begin{bmatrix} 10 \\ -5 \\ 8 \end{bmatrix} \]

\textbf{Solution:}
Set up the augmented matrix:
\[ \begin{bmatrix} 1 & -6 & -5 & 10 \\ 0 & 3 & 2 & -5 \\ 2 & -12 & 9 & 8 \end{bmatrix} \]
Perform row reduction. Eliminate the $2$ in the third row:
\[ \underrightarrow{R_3 \leftarrow R_3 - 2R_1} \]
Calculate $R_3$: $\begin{bmatrix} 2-2 & -12-(-12) & 9-(-10) & 8-20 \end{bmatrix} = \begin{bmatrix} 0 & 0 & 19 & -12 \end{bmatrix}$
The new matrix is:
\[ \begin{bmatrix} 1 & -6 & -5 & 10 \\ 0 & 3 & 2 & -5 \\ 0 & 0 & 19 & -12 \end{bmatrix} \]
The matrix is in echelon form. The pivots are the $1$ in col 1, the $3$ in col 2, and the $19$ in col 3. There is no pivot in the augmented column (col 4). Since the system is consistent, $\mathbf{b}$ is a linear combination of the columns of $A$.
\\
\textbf{Answer:} B. Vector $\mathbf{b}$ is a linear combination... The pivots... are in the first entry in the first column, the second entry in the second column, and the third entry in the third column.


\subsection*{Problem 1.3.17}
\textbf{Question:} Let $\mathbf{a}_1 = \begin{bmatrix} 1 \\ 5 \\ -1 \end{bmatrix}, \mathbf{a}_2 = \begin{bmatrix} -6 \\ -23 \\ 2 \end{bmatrix}$, and $\mathbf{b} = \begin{bmatrix} 3 \\ 1 \\ h \end{bmatrix}$. For what value(s) of $h$ is $\mathbf{b}$ in the plane spanned by $\mathbf{a}_1$ and $\mathbf{a}_2$?

\textbf{Solution:}
Being in the plane spanned by $\mathbf{a}_1$ and $\mathbf{a}_2$ means $\mathbf{b}$ is a linear combination of them. Thus, the augmented matrix $[\mathbf{a}_1 \ \mathbf{a}_2 \ \mathbf{b}]$ must be consistent.
\[ \begin{bmatrix} 1 & -6 & 3 \\ 5 & -23 & 1 \\ -1 & 2 & h \end{bmatrix} \]
Row reduction:
\[ \underrightarrow{R_2 \leftarrow R_2 - 5R_1, \ R_3 \leftarrow R_3 + R_1} \begin{bmatrix} 1 & -6 & 3 \\ 0 & 7 & -14 \\ 0 & -4 & h+3 \end{bmatrix} \]
To avoid fractions, observe that in row 2, $7x_2 = -14 \implies x_2 = -2$.
Substitute $x_2 = -2$ into row 3's equation:
\[ -4x_2 = h + 3 \]
\[ -4(-2) = h + 3 \implies 8 = h + 3 \implies h = 5 \]
Alternatively, by strict row reduction, $R_3 \leftarrow R_3 + \frac{4}{7}R_2$ yields a bottom row of $[0 \ 0 \ h+3-8]$. For consistency, $h - 5 = 0 \implies h = 5$.
\\
\textbf{Answer:} 5


\subsection*{Problem 1.3.19}
\textbf{Question:} Give a geometric description of $\text{Span}\{\mathbf{v}_1, \mathbf{v}_2\}$ for the vectors $\mathbf{v}_1 = \begin{bmatrix} 6 \\ 3 \\ -15 \end{bmatrix}$ and $\mathbf{v}_2 = \begin{bmatrix} 8 \\ 4 \\ -20 \end{bmatrix}$.

\textbf{Solution:}
Normally, the span of two vectors in $\mathbb{R}^3$ forms a plane. However, let's examine the relationship between these two vectors.
Notice that $\mathbf{v}_2 = \frac{4}{3} \mathbf{v}_1$.
Because $\mathbf{v}_2$ is a scalar multiple of $\mathbf{v}_1$, they lie on the exact same line passing through the origin. Therefore, their combined span is no greater than the span of $\mathbf{v}_1$ alone. The span is just that single line.
\\
\textbf{Answer:} C. $\text{Span}\{\mathbf{v}_1, \mathbf{v}_2\}$ is the set of points on the line through $\mathbf{v}_1$ and $\mathbf{0}$.


\subsection*{Problem 1.3.22}
\textbf{Question:} Construct a $3 \times 3$ matrix $A$, with nonzero entries, and a vector $\mathbf{b}$ in $\mathbb{R}^3$ such that $\mathbf{b}$ is not in the set spanned by the columns of $A$.

\textbf{Solution:}
To ensure $\mathbf{b}$ is NOT in the span, we need an inconsistent system. The easiest way to restrict the span of a $3 \times 3$ matrix is to make all its columns identical (or multiples of each other). This restricts the span to a single 1D line in 3D space, leaving almost the entire 3D space outside the span.
Look at Option C:
\[ A = \begin{bmatrix} 1 & 1 & 1 \\ 2 & 2 & 2 \\ 3 & 3 & 3 \end{bmatrix}, \ \mathbf{b} = \begin{bmatrix} 4 \\ 5 \\ 6 \end{bmatrix} \]
All columns of $A$ are multiples of $\begin{bmatrix} 1 \\ 2 \\ 3 \end{bmatrix}$. The vector $\mathbf{b} = \begin{bmatrix} 4 \\ 5 \\ 6 \end{bmatrix}$ is not a multiple of $\begin{bmatrix} 1 \\ 2 \\ 3 \end{bmatrix}$ (e.g., $1 \times 4 = 4$, but $2 \times 4 \neq 5$). Therefore, it cannot be formed by them.
\\
\textbf{Answer:} C.


\subsection*{Problem 1.3.29}
\textbf{Question:} Determine whether the statement below is true or false. 
\textit{The solution set of the linear system whose augmented matrix is $[\mathbf{a}_1 \ \mathbf{a}_2 \ \mathbf{a}_3 \ \mathbf{b}]$ is the same as the solution set of the equation $x_1\mathbf{a}_1 + x_2\mathbf{a}_2 + x_3\mathbf{a}_3 = \mathbf{b}$.}

\textbf{Solution:}
This is the fundamental equivalence theorem of Section 1.3. A vector equation $x_1\mathbf{a}_1 + \dots + x_n\mathbf{a}_n = \mathbf{b}$ possesses the exact same solution weights ($x_1, \dots, x_n$) as the system of linear equations represented by the augmented matrix containing those vectors as columns.
\\
\textbf{Answer:} C. True. The augmented matrix for $x_1\mathbf{a}_1 + x_2\mathbf{a}_2 + x_3\mathbf{a}_3 = \mathbf{b}$ is $[\mathbf{a}_1 \ \mathbf{a}_2 \ \mathbf{a}_3 \ \mathbf{b}]$.


\subsection*{Problem 1.3.39}
\textbf{Question:} A thin triangular plate of uniform density... has vertices at $\mathbf{v}_1=(1,1), \mathbf{v}_2=(7,1), \mathbf{v}_3=(3,7)$, and the mass is $3g$.
\textbf{a.} Find the $(x,y)$-coordinates of the center of mass.
\textbf{b.} Determine how to distribute an additional mass of $6g$ at the vertices to move the balance point to $(3,3)$. Let $w_1, w_2, w_3$ denote the added masses.

\textbf{Solution:}
\textbf{Part a:} The center of mass of 3 equal $1g$ masses is simply the average of their coordinates.
\[ \text{CM} = \frac{1}{3}(\mathbf{v}_1 + \mathbf{v}_2 + \mathbf{v}_3) = \frac{1}{3}\left( \begin{bmatrix} 1 \\ 1 \end{bmatrix} + \begin{bmatrix} 7 \\ 1 \end{bmatrix} + \begin{bmatrix} 3 \\ 7 \end{bmatrix} \right) = \frac{1}{3}\begin{bmatrix} 11 \\ 9 \end{bmatrix} = \begin{bmatrix} 11/3 \\ 3 \end{bmatrix} \]
\textbf{Part b:} We add masses $w_1, w_2, w_3$ to the vertices such that $w_1 + w_2 + w_3 = 6$. The total mass of the system becomes $3g + 6g = 9g$.
The new mass at $\mathbf{v}_1$ is $(1 + w_1)$, at $\mathbf{v}_2$ is $(1 + w_2)$, and at $\mathbf{v}_3$ is $(1 + w_3)$.
The new center of mass equation is:
\[ \frac{1}{9}\left[ (1+w_1)\begin{bmatrix} 1 \\ 1 \end{bmatrix} + (1+w_2)\begin{bmatrix} 7 \\ 1 \end{bmatrix} + (1+w_3)\begin{bmatrix} 3 \\ 7 \end{bmatrix} \right] = \begin{bmatrix} 3 \\ 3 \end{bmatrix} \]
Multiply by 9 to clear the fraction:
\[ \begin{bmatrix} 1+w_1 + 7+7w_2 + 3+3w_3 \\ 1+w_1 + 1+w_2 + 7+7w_3 \end{bmatrix} = \begin{bmatrix} 27 \\ 27 \end{bmatrix} \]
This yields two equations. Combined with the total mass equation, we have a system of three linear equations:
1. $w_1 + w_2 + w_3 = 6$
2. $w_1 + 7w_2 + 3w_3 + 11 = 27 \implies w_1 + 7w_2 + 3w_3 = 16$
3. $w_1 + w_2 + 7w_3 + 9 = 27 \implies w_1 + w_2 + 7w_3 = 18$

Set up the augmented matrix:
\[ \begin{bmatrix} 1 & 1 & 1 & 6 \\ 1 & 7 & 3 & 16 \\ 1 & 1 & 7 & 18 \end{bmatrix} \]
Row reduction:
\[ \underrightarrow{R_2 \leftarrow R_2 - R_1, \ R_3 \leftarrow R_3 - R_1} \begin{bmatrix} 1 & 1 & 1 & 6 \\ 0 & 6 & 2 & 10 \\ 0 & 0 & 6 & 12 \end{bmatrix} \]
Solve by back substitution:
$6w_3 = 12 \implies w_3 = 2$
$6w_2 + 2(2) = 10 \implies 6w_2 = 6 \implies w_2 = 1$
$w_1 + 1 + 2 = 6 \implies w_1 = 3$
\\
\textbf{Answer a:} $(11/3, 3)$. 
\textbf{Answer b:} Add 3g at (1,1), add 1g at (7,1), add 2g at (3,7).


\subsection*{Problem 1.3.33}
\textbf{Question:} Let $A = \begin{bmatrix} 1 & 0 & -6 \\ 0 & 2 & -3 \\ -6 & 4 & 3 \end{bmatrix}$ and $\mathbf{b} = \begin{bmatrix} 16 \\ 1 \\ -13 \end{bmatrix}$. Let $W = \text{Span}\{\mathbf{a}_1, \mathbf{a}_2, \mathbf{a}_3\}$.
\textbf{a.} Is $\mathbf{b}$ in $\{\mathbf{a}_1, \mathbf{a}_2, \mathbf{a}_3\}$? How many vectors?
\textbf{b.} Is $\mathbf{b}$ in $W$? How many vectors in $W$?
\textbf{c.} Show that $\mathbf{a}_2$ is in $W$.

\textbf{Solution:}
\textbf{Part a:} The set $\{\mathbf{a}_1, \mathbf{a}_2, \mathbf{a}_3\}$ is a discrete list containing exactly 3 vectors. Since $\mathbf{b}$ is not identically equal to any of those three columns, it is not in the set.
\textbf{Part b:} $W$ is the \textit{Span} of the vectors, meaning the set of all possible linear combinations. To check if $\mathbf{b}$ is in $W$, we row reduce $[A \ \mathbf{b}]$:
\[ \begin{bmatrix} 1 & 0 & -6 & 16 \\ 0 & 2 & -3 & 1 \\ -6 & 4 & 3 & -13 \end{bmatrix} \]
\[ \underrightarrow{R_3 \leftarrow R_3 + 6R_1} \begin{bmatrix} 1 & 0 & -6 & 16 \\ 0 & 2 & -3 & 1 \\ 0 & 4 & -33 & 83 \end{bmatrix} \]
\[ \underrightarrow{R_3 \leftarrow R_3 - 2R_2} \begin{bmatrix} 1 & 0 & -6 & 16 \\ 0 & 2 & -3 & 1 \\ 0 & 0 & -27 & 81 \end{bmatrix} \]
The system is consistent. Solving it:
$-27x_3 = 81 \implies x_3 = -3$
$2x_2 - 3(-3) = 1 \implies 2x_2 + 9 = 1 \implies 2x_2 = -8 \implies x_2 = -4$
$x_1 - 6(-3) = 16 \implies x_1 + 18 = 16 \implies x_1 = -2$
So $\mathbf{b} = -2\mathbf{a}_1 - 4\mathbf{a}_2 - 3\mathbf{a}_3$. Because $W$ is a span, it contains an infinite number of linear combinations (vectors).
\textbf{Part c:} Any vector is in its own span because you can weight it by 1 and all others by 0.
\\
\textbf{Answers:}
\textbf{Part a:} B. No, $\mathbf{b}$ is not in $\{\mathbf{a}_1, \mathbf{a}_2, \mathbf{a}_3\}$ since $\mathbf{b}$ is not equal to $\mathbf{a}_1, \mathbf{a}_2,$ or $\mathbf{a}_3$. There are 3 vectors in the set.
\textbf{Part b:} A. Yes, $\mathbf{b}$ is in $W$ ... $\mathbf{b} = (-2)\mathbf{a}_1 + (-4)\mathbf{a}_2 + (-3)\mathbf{a}_3$. There are infinitely many vectors in $W$.
\textbf{Part c:} The vector $\mathbf{a}_2$ can be written as $(0)\mathbf{a}_1 + (1)\mathbf{a}_2 + (0)\mathbf{a}_3$.

\newpage
\section{Part 3: In-Depth Analysis of Problems and Techniques}

\subsection{A) Problem Types and General Approach}
\textbf{1. Vector Arithmetic and Translation (Problems 1.3.1, 1.3.5, 1.3.6, 1.3.9):}
\begin{itemize}
    \item \textit{Characteristics:} Basic operations. Translating back and forth between systems of equations and vector equations.
    \item \textit{Approach:} Apply definitions directly. To convert a system to a vector equation, group all $x_1$ coefficients into a vector, all $x_2$ into a vector, etc., and factor out the variables as scalar weights.
\end{itemize}

\textbf{2. Membership in a Span (Problems 1.3.11, 1.3.13, 1.3.14, 1.3.17):}
\begin{itemize}
    \item \textit{Characteristics:} Asking "Is vector $\mathbf{b}$ a linear combination of..." or "Is $\mathbf{b}$ in the plane spanned by..."
    \item \textit{Approach:} Immediately translate the problem into an augmented matrix $[\mathbf{a}_1 \ \dots \ \mathbf{a}_n \ \mathbf{b}]$. Perform row reduction. If the system is consistent (no $[0 \dots 0 \mid b \neq 0]$ row), the answer is Yes. If it is inconsistent, the answer is No.
\end{itemize}

\textbf{3. Geometric Interpretation (Problems 1.3.19, 1.3.22):}
\begin{itemize}
    \item \textit{Characteristics:} Connecting algebraic spans to geometric shapes (lines, planes).
    \item \textit{Approach:} Inspect the vectors before calculating. If two vectors are scalar multiples of one another, they are collinear, and their span is just a 1D line, not a 2D plane. 
\end{itemize}

\textbf{4. Conceptual Distinctions (Problem 1.3.33):}
\begin{itemize}
    \item \textit{Characteristics:} Testing the difference between a Set $\{\mathbf{v}_1, \mathbf{v}_2\}$ and a Span $\text{Span}\{\mathbf{v}_1, \mathbf{v}_2\}$.
    \item \textit{Approach:} Remember that a Set is a rigid bucket containing exactly $n$ specific items. A Span is a manufacturing plant that can produce infinitely many items by mixing the original $n$ ingredients.
\end{itemize}

\subsection{B) Key Algebraic and Calculus Manipulations}
\begin{itemize}
    \item \textbf{The Linear Combination to Augmented Matrix Pipeline:} The most critical algebraic technique is realizing that writing out $x_1\mathbf{a}_1 + x_2\mathbf{a}_2 = \mathbf{b}$ is algebraically redundant. You can immediately extract the columns $\mathbf{a}_1$ and $\mathbf{a}_2$, place them in a matrix, append $\mathbf{b}$, and run Gaussian elimination. 
    \item \textbf{Parameter Isolation:} In Problem 1.3.17, a variable $h$ was buried in the augmented column. The trick is to row reduce the coefficient side strictly to echelon form, letting the $h$ parameter get modified algebraically (e.g., $h+3$). Then, force consistency by analyzing the resulting row.
\end{itemize}

\newpage
\section{Part 4: "Cheatsheet" and Tips for Success}

\subsection{Key Definitions and Formulas}
\begin{itemize}
    \item \textbf{Vector Equation Form:} $x_1\mathbf{a}_1 + x_2\mathbf{a}_2 + \dots + x_n\mathbf{a}_n = \mathbf{b}$
    \item \textbf{Matrix Equation Form:} $A\mathbf{x} = \mathbf{b}$ (where $A$ has columns $\mathbf{a}_1 \dots \mathbf{a}_n$)
    \item \textbf{The Golden Rule of Section 1.3:} A vector equation has the \textbf{same solution set} as the linear system whose augmented matrix is $[\mathbf{a}_1 \ \mathbf{a}_2 \ \dots \ \mathbf{a}_n \ \mathbf{b}]$.
\end{itemize}

\subsection{Shortcuts and Cheats}
\begin{itemize}
    \item \textbf{The Collinearity Check:} Before doing row operations to find the span of two vectors, check if one is a multiple of the other (e.g., $\begin{bmatrix} 2 \\ 4 \end{bmatrix}$ and $\begin{bmatrix} -6 \\ -12 \end{bmatrix}$). If they are, the span is a line. If they are not, the span is a plane.
    \item \textbf{The Zero Vector Trick:} The zero vector $\mathbf{0}$ is \textit{always} in the span of any set of vectors. Just set all weights $x_i = 0$.
    \item \textbf{Self-Containment:} Every vector $\mathbf{v}_i$ in a set is automatically inside its own Span. Set $x_i = 1$ and all other weights to $0$.
\end{itemize}

\subsection{Common Pitfalls and Traps}
\begin{itemize}
    \item \textbf{Confusing "Set" with "Span":} If asked "Is $\mathbf{b}$ in $\{\mathbf{v}_1, \mathbf{v}_2\}$?", only answer Yes if $\mathbf{b}$ looks exactly like $\mathbf{v}_1$ or $\mathbf{v}_2$. Do not row reduce.
    \item \textbf{Dimension Mismatch Error:} You cannot add vectors of different sizes (e.g., adding a vector in $\mathbb{R}^2$ to a vector in $\mathbb{R}^3$ is undefined).
    \item \textbf{Assuming a matrix spans the whole space:} Just because a matrix is $3 \times 3$ does not mean its columns span all of $\mathbb{R}^3$. If the columns are linearly dependent (e.g., all multiples of each other like in Problem 1.3.22), they might only span a plane or a line.
\end{itemize}

\newpage
\section{Part 5: Conceptual Synthesis and The "Big Picture"}

\subsection{A) Thematic Connections}
The core theme of Section 1.3 is \textbf{geometric synthesis}. In earlier math classes, algebra (solving equations) and geometry (drawing lines) were often treated as separate worlds. Linear algebra merges them. A system of equations is no longer just a puzzle to find variables; it is a question of whether you can reach a specific destination (vector $\mathbf{b}$) by traveling along specific paths (vectors $\mathbf{a}_1, \dots, \mathbf{a}_n$).

\subsection{B) Forward and Backward Links}
\begin{itemize}
    \item \textbf{Backward Link:} This section repackages the raw computation of row reduction from 1.1 and 1.2 into a more useful format. Instead of dealing with abstract variables $x, y, z$, we are dealing with discrete packets of data (vectors).
    \item \textbf{Forward Link:} In the very next section (1.4), we formalize the notation $x_1\mathbf{a}_1 + \dots + x_n\mathbf{a}_n$ into the matrix-vector product $A\mathbf{x}$. Furthermore, the concept of "Span" is the bedrock of understanding Vector Spaces and Subspaces in Chapter 4. Understanding linear combinations is required to understand \textbf{linear independence} in Section 1.7.
\end{itemize}

\newpage
\section{Part 6: Real-World Application and Modeling}

\subsection{A) Concrete Scenarios}
\begin{enumerate}
    \item \textbf{Computer Graphics (Color Spaces):} In digital screens, every pixel's color is a linear combination of three primary color vectors: Red, Green, and Blue. Changing the scalar weights (from 0 to 255) generates the entire Span of visible colors on the monitor. If a specific color (vector $\mathbf{b}$) is out of gamut, it means it is not in the span of the monitor's specific RGB vectors.
    \item \textbf{Dietary/Resource Optimization:} A dietician is trying to construct a meal plan that meets exact macronutrient goals (Protein, Carbs, Fat). They have three foods (e.g., Chicken, Rice, Avocado). Each food is a vector in $\mathbb{R}^3$ representing its macronutrient profile per 100g. Finding the correct portion sizes is equivalent to solving a vector equation to reach the target "goal vector."
\end{enumerate}

\subsection{B) Model Problem Setup: Color Space Mixing}
Imagine a lighting engineer has three spotlights. 
Spotlight 1 emits color $\mathbf{c}_1 = \begin{bmatrix} 70 \\ 20 \\ 10 \end{bmatrix}$ (mostly red).
Spotlight 2 emits color $\mathbf{c}_2 = \begin{bmatrix} 10 \\ 60 \\ 20 \end{bmatrix}$ (mostly green).
Spotlight 3 emits color $\mathbf{c}_3 = \begin{bmatrix} 10 \\ 10 \\ 80 \end{bmatrix}$ (mostly blue).
The engineer wants to produce a specific shade of white light given by the target vector $\mathbf{w} = \begin{bmatrix} 150 \\ 150 \\ 150 \end{bmatrix}$.

The engineer must find the intensity settings (scalars $x_1, x_2, x_3$) for each spotlight.
The mathematical model is the vector equation:
\[ x_1 \begin{bmatrix} 70 \\ 20 \\ 10 \end{bmatrix} + x_2 \begin{bmatrix} 10 \\ 60 \\ 20 \end{bmatrix} + x_3 \begin{bmatrix} 10 \\ 10 \\ 80 \end{bmatrix} = \begin{bmatrix} 150 \\ 150 \\ 150 \end{bmatrix} \]
To solve this, they would place the vectors into an augmented matrix and row reduce. If the solution requires a negative weight (e.g., $x_1 = -2$), it means the color is physically impossible to produce with those lights, because you cannot output "negative" light.

\newpage
\section{Part 7: Common Variations and Untested Concepts}

A vital concept taught in 1.3 that was not directly tested in your homework is proving whether a set of vectors spans the \textbf{entire} space $\mathbb{R}^n$, rather than just checking if a single vector $\mathbf{b}$ is in the span.

\subsection{Variation: Spanning all of $\mathbb{R}^n$}
\textbf{Concept:} A set of vectors $\{\mathbf{v}_1, \dots, \mathbf{v}_p\}$ spans all of $\mathbb{R}^n$ if every single possible vector $\mathbf{b}$ in $\mathbb{R}^n$ can be written as a linear combination of them.
This occurs if and only if the matrix $A = [\mathbf{v}_1 \dots \mathbf{v}_p]$ has a pivot position in \textbf{every row}.

\textbf{Example Problem:} Do the vectors $\mathbf{v}_1 = \begin{bmatrix} 1 \\ 0 \\ -2 \end{bmatrix}, \mathbf{v}_2 = \begin{bmatrix} 3 \\ 1 \\ -4 \end{bmatrix}, \mathbf{v}_3 = \begin{bmatrix} -2 \\ -1 \\ 2 \end{bmatrix}$ span $\mathbb{R}^3$?

\textbf{Setup and Solution:}
We place the vectors into a matrix and row reduce to count the pivots.
\[ A = \begin{bmatrix} 1 & 3 & -2 \\ 0 & 1 & -1 \\ -2 & -4 & 2 \end{bmatrix} \]
\[ \underrightarrow{R_3 \leftarrow R_3 + 2R_1} \begin{bmatrix} 1 & 3 & -2 \\ 0 & 1 & -1 \\ 0 & 2 & -2 \end{bmatrix} \]
\[ \underrightarrow{R_3 \leftarrow R_3 - 2R_2} \begin{bmatrix} 1 & 3 & -2 \\ 0 & 1 & -1 \\ 0 & 0 & 0 \end{bmatrix} \]
The matrix has pivots in Row 1 and Row 2, but \textbf{no pivot in Row 3}. 
Because there is a row without a pivot, there will exist some target vectors $\mathbf{b}$ that result in an inconsistent system (specifically, vectors that don't satisfy the condition generated by the zero row). 
Therefore, the vectors do \textbf{not} span all of $\mathbb{R}^3$. They only span a 2D plane within $\mathbb{R}^3$.

\newpage
\section{Part 8: Advanced Diagnostic Testing: "Find the Flaw"}

Below are 5 problems with complete, seemingly correct solutions. However, each contains a single fatal flaw. Identify the error and provide the correct calculation.

\hrule
\vspace{0.5cm}

\textbf{Flawed Problem 1:} Compute the linear combination $3\mathbf{u} - \mathbf{v}$ where $\mathbf{u} = \begin{bmatrix} 1 \\ 2 \\ 0 \end{bmatrix}$ and $\mathbf{v} = \begin{bmatrix} 4 \\ -1 \end{bmatrix}$.
\textit{Flawed Solution:} 
$3\mathbf{u} = \begin{bmatrix} 3 \\ 6 \\ 0 \end{bmatrix}$. Subtracting $\mathbf{v}$ yields $\begin{bmatrix} 3 - 4 \\ 6 - (-1) \\ 0 - 0 \end{bmatrix} = \begin{bmatrix} -1 \\ 7 \\ 0 \end{bmatrix}$.

\vspace{0.3cm}
\textbf{Flawed Problem 2:} Determine if $\mathbf{b} = \begin{bmatrix} 2 \\ -1 \end{bmatrix}$ is in the set spanned by $\mathbf{v}_1 = \begin{bmatrix} 1 \\ 0 \end{bmatrix}$ and $\mathbf{v}_2 = \begin{bmatrix} 0 \\ 1 \end{bmatrix}$.
\textit{Flawed Solution:} 
Because $\mathbf{b}$ is not equal to $\mathbf{v}_1$ and $\mathbf{b}$ is not equal to $\mathbf{v}_2$, it is not contained within the span of the two vectors. The answer is No.

\vspace{0.3cm}
\textbf{Flawed Problem 3:} Describe the geometric span of $\mathbf{v}_1 = \begin{bmatrix} 1 \\ 2 \\ 3 \end{bmatrix}$ and $\mathbf{v}_2 = \begin{bmatrix} 2 \\ 4 \\ 6 \end{bmatrix}$.
\textit{Flawed Solution:} 
Because there are two non-zero vectors in $\mathbb{R}^3$, they define two distinct directions. Therefore, the span of these two vectors geometrically forms a 2-dimensional plane passing through the origin in $\mathbb{R}^3$.

\vspace{0.3cm}
\textbf{Flawed Problem 4:} Determine if $\mathbf{b}$ is a linear combination of the columns of $A$.
$A = \begin{bmatrix} 1 & 2 \\ 0 & 3 \end{bmatrix}, \ \mathbf{b} = \begin{bmatrix} 4 \\ 6 \end{bmatrix}$.
\textit{Flawed Solution:} 
Augmented matrix: $\begin{bmatrix} 1 & 2 & 4 \\ 0 & 3 & 6 \end{bmatrix}$.
Divide Row 2 by 3: $\begin{bmatrix} 1 & 2 & 4 \\ 0 & 1 & 2 \end{bmatrix}$.
The matrix is in echelon form. The pivots are the 1 in column 1, and the 2 in column 3. Because there is a pivot in the rightmost column, the system is inconsistent, and $\mathbf{b}$ is not a linear combination.

\vspace{0.3cm}
\textbf{Flawed Problem 5:} Write the system of equations equivalent to $x_1 \begin{bmatrix} 1 \\ 2 \end{bmatrix} - x_2 \begin{bmatrix} 3 \\ 4 \end{bmatrix} = \begin{bmatrix} 5 \\ 6 \end{bmatrix}$.
\textit{Flawed Solution:} 
Applying the scalars yields:
$1x_1 - 3x_2 = 5$
$2x_1 + 4x_2 = 6$
(The second equation uses $+4x_2$ because we are combining the vectors).

\end{document}